\section{Method}
\label{sec:method}

\subsection{\ours{}: Two Cosine Schedules, One Ensemble}

Let $f_\theta$ be a neural network trained with a cosine-annealed
learning-rate schedule~\cite{loshchilov2017sgdr} starting at
$\mathrm{lr}_{\mathrm{init}}$ and decaying to $\mathrm{lr}_{\mathrm{min}}$
over $T_\mathrm{max}$ epochs.
Standard deep ensembling trains $N$ independent copies of $f_\theta$ with
$N$ different random seeds at the same $T_\mathrm{max}$ and averages softmax
outputs at test time.

\ours{} trains $2N$ models: $N$ seeds with $T_\mathrm{max}=T_A$ and $N$
seeds with $T_\mathrm{max}=T_B$ (with $T_B > T_A$, e.g., $T_B = 1.5 T_A$).
The final ensemble averages softmax outputs from all $2N$ models:
\[
\bar p(x) = \frac{1}{2N}\sum_{i=1}^{2N} \mathrm{softmax}\!\bigl(f_{\theta_i}(x)\bigr).
\]

There is no change to the training objective, loss function, random seeds,
augmentation pipeline, optimizer, or any other component. The only change
is a single integer (the endpoint of the cosine schedule) differing across
the two groups.

\subsection{Why Two Endpoints?}

\paragraph{The ``low-LR exploration'' phase.}
A cosine schedule $\mathrm{lr}(t) = \mathrm{lr}_{\mathrm{min}} +
\frac{1}{2}(\mathrm{lr}_{\mathrm{init}} - \mathrm{lr}_{\mathrm{min}})(1 + \cos(\pi t / T_{\max}))$
at $t=100$ with $T_\mathrm{max}=100$ produces
$\mathrm{lr}=\mathrm{lr}_{\mathrm{min}} \approx 10^{-6}$. The same schedule at
$t=100$ with $T_\mathrm{max}=150$ produces
$\mathrm{lr} = \mathrm{lr}_{\mathrm{init}} \cdot (1 + \cos(2\pi/3))/2 = 0.25 \cdot \mathrm{lr}_{\mathrm{init}} \approx 2.5 \times 10^{-4}$.
That is a 250$\times$ difference in instantaneous LR.

So a model training for 150 epochs, at epoch 100, is still in the
``moderate-LR'' regime and continues to take meaningful gradient steps
through the loss landscape. Over the remaining 50 epochs, the LR slowly
decays to $\mathrm{lr}_{\mathrm{min}}$. During this 50-epoch tail, the
model \emph{stochastically explores} the loss basin it is in: each SGD
step perturbs the weights, but the decreasing LR eventually settles it
into some local minimum of that basin.

\paragraph{Empirical finding.}
We show in \cref{sec:experiments} that:
(i) the 150-epoch and 100-epoch runs reach \emph{identical} mean
validation accuracy (difference $<\!0.001$);
(ii) but they produce different consensus predictions on 13\% of test samples;
(iii) on those disagreement samples, the mixed ensemble is 10 percentage
points more accurate than either group's consensus;
(iv) multiplying 10\,pp $\times$ 13\% gives $+1.3$~pp total gain, closely
matching the observed $+1.5$~pp between the same-length ensemble and the
mixed-length ensemble.

\subsection{Choosing the Endpoint Ratio: Why $T_B = 1.5\,T_A$?}
\label{sec:why_150}

The endpoint ratio $\rho = T_B / T_A$ is the single user-visible
hyperparameter of \ours. Two considerations constrain it. (i)~$\rho$
must be large enough that group $B$'s tail phase visits a genuinely
different local minimum: the PCA analysis in \cref{sec:basin} confirms
that $\rho = 1.5$ produces a between-group centroid separation
$3.25\times$ larger than the within-group pairwise distance. (ii)~$\rho$
must be small enough that group $B$ does not overfit or drift into a
low-accuracy region. A direct sweep on FACED EMOD d6 with
$T_A=100$ and $T_B \in \{110, 120, 130, 140, 150, 200\}$
(\cref{sec:endpoint_sweep}) shows that the gain over the single-length
$T_A$ ensemble is positive and largest at $\rho \approx 1.5$. We therefore
recommend $\rho \in [1.4, 1.6]$ as the operating regime and use $\rho = 1.5$
throughout the rest of the paper; the ratio is task-robust rather than
task-tuned.

\subsection{Theoretical Prediction of Gain Magnitude}

Let $p_A$ be the accuracy of a single seed trained with $T_\mathrm{max}=T_A$
and let $p_B$ be the same for $T_B$. Let $q$ be the fraction of test samples
on which the 100-epoch and 150-epoch \emph{group consensuses} disagree, and
let $g$ be the accuracy gain on disagreement samples when using the mixed
ensemble over either single-length ensemble. The predicted total gain is
\[
\Delta_{\mathrm{mixed}} \approx q \cdot g.
\]

We verify this equation empirically: on FACED d6 ensembles, $q = 0.131$,
$g = 0.103$, product $= 0.0135$, observed gain $= 0.015$.

\subsection{Gain-vs-Accuracy Scaling Law}

We also conjecture and verify empirically that the gain scales inversely
with individual accuracy. This follows from a simple argument: if
individual models are already perfect, there are no disagreement samples
and no gain from ensembling (mixed or otherwise). If individual models are
near-random, many disagreement samples are also near-random and ensembling
helps less. The sweet spot is individual accuracy in the 0.3--0.8 range,
where many disagreement samples are still decodable by the richer
ensemble distribution.

\subsection{Complexity}

\paragraph{Compute.} \ours{} requires two training runs per seed, one with
$T_\mathrm{max}=T_A$ and one with $T_\mathrm{max}=T_B$. Total compute is
$\sim\!T_A + T_B$ epochs per seed, vs.\ $\sim\!2T_A$ for a standard $2N$-seed
deep ensemble with the same total ensemble size. The extra cost is
$(T_B - T_A) \cdot N$ training epochs (e.g., $50 \times 5 = 250$ epochs
for our $T_A=100, T_B=150, N=5$ setup on CIFAR-10). This is a modest
overhead for the $+1.5$~pp gain beyond standard deep ensembles.

\paragraph{Inference.} Identical to standard deep ensembles: $2N$ forward
passes per input, softmax averaged.

\paragraph{Storage.} Identical to standard deep ensembles: $2N$ model
checkpoints.
